\documentclass[12pt, a4paper]{book}

\usepackage[T1]{fontenc}
\usepackage[utf8]{inputenc}
\usepackage[charter, cal=cmcal]{mathdesign}
\usepackage[scaled=0.92]{helvet}
\usepackage{geometry}
\geometry{a4paper, top=30mm, bottom=30mm, inner=32mm, outer=25mm, marginparwidth=18mm}

\usepackage{amsmath, amsthm}
\usepackage[version=4]{mhchem}
\usepackage{graphicx}
\usepackage{xcolor}
\usepackage{tikz}
\usetikzlibrary{arrows.meta, shapes.geometric, positioning, calc, backgrounds}
\usepackage{pgfplots}
\pgfplotsset{compat=1.18}
\usepackage{tcolorbox}
\tcbuselibrary{skins, breakable, theorems}
\usepackage{titlesec}
\usepackage{fancyhdr}
\usepackage{enumitem}
\usepackage{booktabs}
\usepackage{array}
\usepackage{parskip}
\usepackage{siunitx}
\usepackage[colorlinks=true, linkcolor=BrickRed,
            urlcolor=NavyBlue, citecolor=ForestGreen]{hyperref}

% ── Colours (Journey 7 — forest green / amber) ────────────────────────────────
\definecolor{JourneyColor}{HTML}{065F46}
\definecolor{ModuleColor}{HTML}{059669}
\definecolor{PuzzleColor}{HTML}{B45309}
\definecolor{DiscoveryColor}{HTML}{D97706}
\definecolor{WarningColor}{HTML}{991B1B}
\definecolor{ProblemColor}{HTML}{6B21A8}
\definecolor{ReflectColor}{HTML}{0D9488}
\definecolor{HistoryColor}{HTML}{78350F}

% ── Boxes ─────────────────────────────────────────────────────────────────────
\newtcolorbox{puzzlebox}[1][]{enhanced,breakable,
  colback=PuzzleColor!6,colframe=PuzzleColor,fonttitle=\bfseries\sffamily,
  title={The Puzzle},attach boxed title to top left={yshift=-2mm,xshift=6mm},
  boxed title style={colback=PuzzleColor},borderline west={3pt}{0pt}{PuzzleColor},#1}

\newtcolorbox{discoverybox}[1][]{enhanced,breakable,
  colback=DiscoveryColor!8,colframe=DiscoveryColor,fonttitle=\bfseries\sffamily,
  title={Discovery},attach boxed title to top left={yshift=-2mm,xshift=6mm},
  boxed title style={colback=DiscoveryColor},borderline west={3pt}{0pt}{DiscoveryColor},#1}

\newtcolorbox{problembox}[1][]{enhanced,breakable,
  colback=ProblemColor!5,colframe=ProblemColor,fonttitle=\bfseries\sffamily,
  title={Problems},attach boxed title to top left={yshift=-2mm,xshift=6mm},
  boxed title style={colback=ProblemColor},borderline west={3pt}{0pt}{ProblemColor},#1}

\newtcolorbox{deepbox}[1][]{enhanced,breakable,
  colback=JourneyColor!5,colframe=JourneyColor,fonttitle=\bfseries\sffamily,
  title={Deep Problems},attach boxed title to top left={yshift=-2mm,xshift=6mm},
  boxed title style={colback=JourneyColor},borderline west={4pt}{0pt}{JourneyColor},#1}

\newtcolorbox{historybox}[1][]{enhanced,breakable,
  colback=HistoryColor!7,colframe=HistoryColor!60,fonttitle=\bfseries\sffamily\small,
  title={Historical Note},attach boxed title to top left={yshift=-2mm,xshift=6mm},
  boxed title style={colback=HistoryColor!80},#1}

\newtcolorbox{reflectionbox}[1][]{enhanced,breakable,
  colback=ReflectColor!6,colframe=ReflectColor,fonttitle=\bfseries\sffamily,
  title={Reflection},attach boxed title to top left={yshift=-2mm,xshift=6mm},
  boxed title style={colback=ReflectColor},borderline west={3pt}{0pt}{ReflectColor},#1}

\newtcolorbox{namedconcept}[2][]{enhanced,
  colback=ModuleColor!6,colframe=ModuleColor,fonttitle=\bfseries\sffamily\large,
  title={#2},attach boxed title to top left={yshift=-2mm,xshift=6mm},
  boxed title style={colback=ModuleColor},#1}

% ── Section formatting ────────────────────────────────────────────────────────
\titleformat{\chapter}[display]
  {\normalfont\sffamily\huge\bfseries\color{JourneyColor}}
  {\textcolor{JourneyColor!60}{\chaptertitlename\ \thechapter}}{10pt}{\Huge}
  [\vspace{2pt}\textcolor{JourneyColor}{\titlerule[1.5pt]}]
\titleformat{\section}
  {\normalfont\sffamily\Large\bfseries\color{ModuleColor}}
  {\thesection}{1em}{}[\textcolor{ModuleColor!40}{\titlerule[0.6pt]}]
\titleformat{\subsection}
  {\normalfont\sffamily\large\bfseries\color{JourneyColor!80}}{\thesubsection}{1em}{}
\titleformat{\subsubsection}
  {\normalfont\sffamily\normalsize\bfseries}{\thesubsubsection}{1em}{}

\pagestyle{fancy}
\fancyhf{}
\fancyhead[LE]{\small\sffamily\textcolor{JourneyColor}{\leftmark}}
\fancyhead[RO]{\small\sffamily\textcolor{ModuleColor}{\rightmark}}
\fancyfoot[C]{\small\thepage}
\renewcommand{\headrulewidth}{0.4pt}

\newcommand{\newterm}[1]{\textit{\textcolor{DiscoveryColor}{#1}}}
\newcommand{\important}[1]{\textcolor{ModuleColor}{\textbf{#1}}}

\begin{document}

\setcounter{chapter}{6}
\chapter[Journey 7 --- Organic Chemistry: The World It Builds]
        {Journey Seven\\[8pt]
         \textcolor{ModuleColor}{\large Organic Chemistry: The World It Builds}}
\setcounter{section}{0}

\section{Module 7.1 --- Chirality in Medicine and Drug Design}
\label{sec:mod7.1}

\subsection*{Motivation}

We mastered in Journey 6 the grammar of organic structures: geometries, stereochemistry, and electron flow.
But abstract molecular geometry has profound real-world consequences when it meets the machinery of life.

The human body is an intrinsically chiral environment.
Our proteins, enzymes, and receptors are composed exclusively of left-handed (L) amino acids.
Consequently, the receptors that transmit signals in our brains and bodies behave like chiral gloves.
An organic drug molecule of a specific handedness will fit into the glove perfectly, triggering a curative response.
Its mirror image, however, may not fit at all, or worse, may fit into a different glove, causing severe toxicity.
This module explores how stereochemistry governs pharmacology, how we model the thermodynamics of drug binding, and how chiral specificity dictates the design of modern medicine.

\subsection{The Levodopa Asymmetry}

We analyze how the chiral blood-brain barrier distinguishes between mirror-image molecules.

\begin{puzzlebox}
\textbf{L-DOPA vs. D-DOPA.}

Parkinson's disease is a neurodegenerative disorder caused by a depletion of dopamine in the brain. Dopamine itself cannot cross the blood-brain barrier from the blood to the brain.
To treat the disease, patients are given a chiral precursor molecule: **L-DOPA** (levodopa). Once inside the brain, an enzyme removes a carboxyl group to convert L-DOPA into active dopamine.
\begin{enumerate}[label=(\alph*), leftmargin=*]
  \item If patients are given the mirror-image molecule, **D-DOPA** (dextrodopa), it is completely ineffective: it cannot cross the blood-brain barrier and is not recognized by the decarboxylase enzyme. Instead, it accumulates in the body and causes bone marrow toxicity.
        Why do the transport proteins in the blood-brain barrier and the brain enzymes distinguish so sharply between these two mirror-image molecules?
  \item Draw a schematic representing a "chiral lock" (a receptor pocket with three specific binding sites: A, B, and C) and a "chiral key" (a molecule with three corresponding groups: 1, 2, and 3).
        Show that one enantiomer can align all three groups with their matching sites simultaneously, whereas its mirror image can align at most two groups at a time.
  \item This geometric model is the **Easson-Stedman hypothesis**. Explain why the difference between three-point binding and two-point binding leads to a large difference in binding free energy ($\Delta G^\circ$).
\end{enumerate}
\end{puzzlebox}

\subsection{Thermodynamics of Receptor Binding}

We model the equilibrium kinetics of a drug binding to a biological receptor.

\begin{explobox}
The binding of a drug ligand (L) to a receptor (R) is a reversible equilibrium:
\[
  \ce{R + L <=> RL},
\]
characterized by the **dissociation constant** $K_d$:
\[
  K_d \;=\; \frac{[\ce{R}][\ce{L}]}{[\ce{RL}]}.
\]
The standard Gibbs free energy of binding is related to $K_d$ by:
\[
  \Delta G^\circ \;=\; R T \ln K_d.
\]
\begin{enumerate}[label=(\alph*), leftmargin=*]
  \item Let the total concentration of receptors be $[\ce{R}]_{\text{tot}} = [\ce{R}] + [\ce{RL}]$. Prove that the fraction of bound receptors $\theta$ is given by:
  \[
    \theta \;=\; \frac{[\ce{L}]}{K_d + [\ce{L}]}.
  \]
  Observe that this has the exact same mathematical form as the Langmuir adsorption isotherm (Module 5.4) and the Michaelis-Menten equation (Module 5.4). Why does this same hyperbolic scaling emerge across physics, chemistry, and biology?
  \item A drug designer measures the binding affinity of two enantiomers at $T = \SI{298}{K}$.
        The active (R)-enantiomer has $K_{d,\text{R}} = \SI{1.0}{\nano M}$ ($1.0 \times 10^{-9}$ M), while the inactive (S)-enantiomer has $K_{d,\text{S}} = \SI{1.0}{\micro M}$ ($1.0 \times 10^{-6}$ M).
        Calculate the difference in binding free energy $\Delta \Delta G^\circ = \Delta G^\circ_{\text{S}} - \Delta G^\circ_{\text{R}}$.
  \item Notice that a minor energy difference of only $\sim \SI{17}{kJ/mol}$ (the energy of a single weak hydrogen bond) leads to a massive, $1000$-fold difference in binding affinity. Why does this make receptor binding highly sensitive to tiny structural deformations?
\end{enumerate}
\end{explobox}

\subsection{Chiral Drugs in the Clinic}

We study how the human body metabolizes different chiral pharmaceuticals.

\begin{explobox}
Because enantiomers can have completely different biological activities, the FDA regulates the stereochemistry of drugs.
\problem **Ibuprofen:** The popular anti-inflammatory ibuprofen is sold as a racemic mixture. The active form is (S)-ibuprofen. The (R)-enantiomer is inactive but is completely harmless because an isomerase enzyme in the human liver converts (R)-ibuprofen into (S)-ibuprofen in vivo.
\problem **Naproxen:** The structurally similar anti-inflammatory naproxen cannot be interconverted. The (S)-enantiomer is the active pain reliever, whereas the (R)-enantiomer is a severe liver toxin. Consequently, naproxen is sold strictly as a pure enantiomer.
\begin{enumerate}[label=(\alph*), leftmargin=*]
  \item Why is it safe to sell ibuprofen as a racemic mixture but highly dangerous for naproxen?
  \item When a pharmaceutical company changes a drug formulation from a racemic mixture to a pure single enantiomer, we call it a **chiral switch**. Explain the clinical advantages of a chiral switch (consider dosage reduction, metabolic load on the liver, and side effects).
\end{enumerate}
\end{explobox}

\newpage

% ─── Section: Discovery ──────────────────────────────────────────────────────
\section{Part 3 --- Discovery}

\begin{discoverybox}
\textbf{What have we built?}

\textbf{First.} The biological environment is homochiral (proteins are built of L-amino acids), making drug receptors structurally asymmetric.

\textbf{Second.} The Easson-Stedman hypothesis models chiral recognition: a minimum of three specific spatial interactions are required for a receptor to distinguish between two enantiomers.

\textbf{Third.} The fraction of receptors occupied by a drug follows the binding equation $\theta = [\ce{L}]/(K_d + [\ce{L}])$. The dissociation constant $K_d$ is exponentially related to the binding free energy $\Delta G^\circ$.

\textbf{Fourth.} Enantiomers of the same drug can behave differently in vivo: one may be curative while the other is inactive, metabolized into the active form (like ibuprofen), or highly toxic (like naproxen or thalidomide).
\end{discoverybox}

\newpage

% ─── Section: Naming Concepts ────────────────────────────────────────────────
\section{Part 4 --- What This Is Called}

\begin{namedconcept}{Pharmacological parameters}
* **Molecular Recognition:** The specific thermodynamic interaction between a receptor and a drug ligand through non-covalent bonds.
* **Dissociation Constant ($K_d$):** The equilibrium constant measuring a ligand's tendency to detach from a receptor. A smaller $K_d$ indicates tighter binding.
* **Pharmacophore:** The specific 3D spatial arrangement of chemical groups (hydrogen bond donors/acceptors, charges, hydrophobic groups) required for biological activity.
* **Chiral Switch:** The replacement of a previously approved racemic drug with a pure single-enantiomer version.
\end{namedconcept}

\newpage

% ─── Section: Solved Exercises ───────────────────────────────────────────────
\section{Part 5 --- Solved Exercises}

\begin{problembox}
\begin{enumerate}[label=\textbf{E\arabic*.}]
  \item \textbf{Calculating Binding Free Energy Difference.}
  
        An anti-viral drug binds to a viral protease. The active (S)-enantiomer has $K_{d,\text{S}} = 2.5 \times 10^{-9}$ M, while the (R)-enantiomer has $K_{d,\text{R}} = 5.0 \times 10^{-7}$ M at \SI{310}{K} (body temperature).
        Calculate the difference in binding free energy $\Delta \Delta G^\circ$ between the two enantiomers.
        
        \textbf{Solution:}
        The binding free energy for each enantiomer is:
        \[
          \Delta G^\circ_{\text{S}} \;=\; R T \ln K_{d,\text{S}}, \qquad \Delta G^\circ_{\text{R}} \;=\; R T \ln K_{d,\text{R}}.
        \]
        The difference in binding free energy is:
        \[
          \Delta \Delta G^\circ \;=\; \Delta G^\circ_{\text{R}} - \Delta G^\circ_{\text{S}} \;=\; R T \ln K_{d,\text{R}} - R T \ln K_{d,\text{S}} \;=\; R T \ln\left( \frac{K_{d,\text{R}}}{K_{d,\text{S}}} \right).
        \]
        Here:
        * $R = 8.314$ J$\cdot$mol$^{-1}$K$^{-1}$.
        * $T = 310$ K.
        * $K_{d,\text{R}}/K_{d,\text{S}} = \frac{5.0 \times 10^{-7}}{2.5 \times 10^{-9}} = 200$.
        
        Substitute these values:
        \[
          \Delta \Delta G^\circ \;=\; \left( 8.314 \text{ J}\cdot\text{mol}^{-1}\text{K}^{-1} \right) \left( 310 \text{ K} \right) \ln(200).
        \]
        Evaluate the natural logarithm: $\ln(200) \approx 5.30$.
        \[
          \Delta \Delta G^\circ \;=\; 2577 \text{ J/mol} \times 5.30 \;=\; 1.37 \times 10^4 \text{ J/mol} \;=\; 13.7 \text{ kJ/mol}.
        \]
        Thus, the difference in binding free energy is \SI{13.7}{kJ/mol}. This tiny energy difference accounts for the $200$-fold selectivity of the viral protease.
        
  \item \textbf{Deriving Receptor fractional Occupancy.}
  
        Derive the expression for the fraction of bound receptors $\theta = \frac{[\ce{RL}]}{[\ce{R}]_{\text{tot}}}$ in terms of free ligand concentration $[\ce{L}]$ and dissociation constant $K_d$.
        
        \textbf{Solution:}
        By definition, the total receptor concentration is:
        \[
          [\ce{R}]_{\text{tot}} \;=\; [\ce{R}] + [\ce{RL}].
        \]
        The fractional occupancy is:
        \[
          \theta \;=\; \frac{[\ce{RL}]}{[\ce{R}]_{\text{tot}}} \;=\; \frac{[\ce{RL}]}{[\ce{R}] + [\ce{RL}]}.
        \]
        From the definition of the dissociation constant:
        \[
          K_d \;=\; \frac{[\ce{R}][\ce{L}]}{[\ce{RL}]} \implies [\ce{R}] \;=\; \frac{K_d [\ce{RL}]}{[\ce{L}]}.
        \]
        Substitute this expression for $[\ce{R}]$ into the occupancy equation:
        \[
          \theta \;=\; \frac{[\ce{RL}]}{\dfrac{K_d [\ce{RL}]}{[\ce{L}]} + [\ce{RL}]}.
        \]
        Divide the numerator and denominator by $[\ce{RL}]$ (assuming $[\ce{RL}] \neq 0$):
        \[
          \theta \;=\; \frac{1}{\dfrac{K_d}{[\ce{L}]} + 1}.
        \]
        Multiply the numerator and denominator by $[\ce{L}]$:
        \[
          \theta \;=\; \frac{[\ce{L}]}{K_d + [\ce{L}]}.
        \]
        This completes the derivation.
\end{enumerate}
\end{problembox}

\newpage

% ─── Section: Problems ───────────────────────────────────────────────────────
\section{Part 6 --- Problems}

\begin{problembox}
\begin{enumerate}[label=\textbf{P\arabic*.}]
  \item \textbf{Easson-Stedman Three-Point Binding.}
  
        Let a receptor active site have three pockets: a hydrogen-bond acceptor, a hydrophobic pocket, and a positively charged electrostatic site.
        A chiral drug molecule has three corresponding functional groups: a hydrogen-bond donor ($-\ce{OH}$), a phenyl ring ($-\ce{C6H5}$), and a carboxylate anion ($-\ce{COO-}$).
        \begin{enumerate}[label=(\alph*)]
          \item Sketch a diagram of the active site and show how the active enantiomer aligns all three groups to bind.
          \item Sketch the mirror-image enantiomer trying to align with the same site. Show that it can align at most two groups at a time.
          \item If each successful interaction releases $\sim \SI{6}{kJ/mol}$ of binding enthalpy, estimate the difference in binding enthalpy $\Delta \Delta H^\circ$ between the two enantiomers.
        \end{enumerate}

  \item \textbf{Enantiomeric Inversion of Ibuprofen.}
  
        In the liver, (R)-ibuprofen is converted to the active (S)-enantiomer through a thioester intermediate. The reaction proceeds by abstracting the hydrogen from the chiral carbon, forming a planar enol intermediate, which is then re-protonated.
        \begin{enumerate}[label=(\alph*)]
          \item Explain why the planar enol intermediate is achiral. (What is the hybridization state of the carbon atom in the enol double bond?)
          \item Why does this planar geometry allow the chiral center to invert configuration upon re-protonation?
        \end{enumerate}
\end{enumerate}
\end{problembox}

\newpage

% ─── Section: Deep Problems ──────────────────────────────────────────────────
\section{Part 7 --- Deep Problems}

\begin{deepbox}
\begin{enumerate}[label=\textbf{D\arabic*.}]
  \item \textbf{Thermodynamics of Drug Rigidification.}
  
        The binding free energy of a drug ligand is $\Delta G^\circ = \Delta H^\circ - T \Delta S^\circ$.
        When a flexible drug molecule binds to a receptor, it must freeze its internal bond rotations, resulting in a large loss of conformational entropy: $\Delta S^\circ_{\text{conformational}} < 0$.
        To design tighter-binding drugs, chemists often "rigidify" the drug molecule, locking it into the active binding conformation by incorporating it into a ring structure.
        \begin{enumerate}[label=(\alph*)]
          \item Explain why locking a drug into its active conformation beforehand increases the binding affinity (decreases $K_d$). Refer to the thermodynamic entropy term $-T\Delta S^\circ$.
          \item Suppose a flexible drug loses $\SI{15}{kJ/mol}$ of conformational entropy upon binding at $T = \SI{298}{K}$.
                If we pre-rigidify this drug into its active conformation without changing its enthalpy of binding, calculate the expected increase in binding affinity (the ratio $K_{d,\text{flexible}}/K_{d,\text{rigid}}$).
        \end{enumerate}

  \item \textbf{Chiral Chromatographic Resolution.}
  
        Because enantiomers have identical physical properties in achiral environments, separating them requires a chiral stationary phase (CSP) in chromatography.
        The column is packed with silica beads coated with a single pure enantiomer of a chiral selector molecule (e.g., cellulose).
        Let the two enantiomers of a racemic drug mixture be $\ce{L_R}$ and $\ce{L_S}$. As they pass down the column, they interact reversibly with the chiral selector ($\ce{CS}$):
        \[
          \ce{CS + L_R <=> CS\cdot L_R}, \qquad K_{d,\text{R}} \;=\; \frac{[\ce{CS}][\ce{L_R}]}{[\ce{CS\cdot L_R}]},
        \]
        \[
          \ce{CS + L_S <=> CS\cdot L_S}, \qquad K_{d,\text{S}} \;=\; \frac{[\ce{CS}][\ce{L_S}]}{[\ce{CS\cdot L_S}]}.
        \]
        \begin{enumerate}[label=(\alph*)]
          \item Show why the difference in binding affinities ($K_{d,\text{R}} \neq K_{d,\text{S}}$) leads to different travel speeds down the column. Which enantiomer elutes (exits) first: the one with the larger or smaller $K_d$?
          \item If $K_{d,\text{R}} = 1.0 \times 10^{-4}$ M and $K_{d,\text{S}} = 2.0 \times 10^{-4}$ M, calculate the ratio of bound fractions $\theta_{\text{R}}/\theta_{\text{S}}$ on the column stationary phase at low ligand concentrations.
        \end{enumerate}
\end{enumerate}
\end{deepbox}

\newpage

% ─── Section: Reflection ─────────────────────────────────────────────────────
\section{Part 8 --- Reflection}

\begin{reflectionbox}
\textbf{Conceptual Reflection}

Consider the following questions:
\begin{enumerate}
  \item **On Biological gloves:** If all L-amino acids in our bodies were suddenly replaced with D-amino acids, would we survive? Would we be able to digest our food or respond to our current drugs?
  \item **On the Easson-Stedman Model:** Why is a minimum of three points of interaction necessary to distinguish chirality in 3D space? Why would a two-point interaction fail to distinguish between mirror images? (Think about the mathematical definition of a plane).
  \item **On Pharmacology Ethics:** If one enantiomer of a drug is active and the other is toxic (like thalidomide), is it ethical to sell the drug as a racemic mixture if a chiral switch is chemically possible? What if the body interconverts them anyway?
\end{enumerate}
\end{reflectionbox}

\medskip
\noindent\textbf{Looking ahead.}
Module~7.2 introduces **Nucleic Acids and DNA Structure**. Having explored the chiral pharmacology of small molecules, we will scale up to the absolute genetic blueprint of life. We will study the specific hydrogen-bonding geometries of base pairs and discover how organic structures encode the information of heredity.

\vspace{1.4cm}
\begin{center}
{\color{JourneyColor}\rule{0.5\linewidth}{0.8pt}}\\[6pt]
{\small\color{JourneyColor}\textit{End of Journey 7, Module 7.1 --- Chirality in Medicine and Drug Design}}\\[4pt]
{\small\color{JourneyColor}Continue to: \textbf{Module 7.2 --- Nucleic Acids and DNA Structure}}\\[6pt]
{\color{JourneyColor}\rule{0.5\linewidth}{0.8pt}}
\end{center}

\end{document}
